\documentclass[11pt,a4paper]{article}

\usepackage[utf8]{inputenc}
\usepackage[T1]{fontenc}
\usepackage[french]{babel}
\usepackage[a4paper,margin=2cm]{geometry}
\usepackage{xcolor}
\usepackage{tikz}
\usetikzlibrary{positioning,arrows.meta}
\usepackage{array}
\usepackage{booktabs}
\usepackage{longtable}
\usepackage{tabularx}
\usepackage{float}
\usepackage{hyperref}
\usepackage{enumitem}
\usepackage{graphicx}
\usepackage{listings}
\usepackage{caption}
\usepackage{titlesec}
\usepackage{fancyhdr}
\usepackage{amsmath}
\usepackage{setspace}
\usepackage{colortbl}

\hypersetup{
    colorlinks=true,
    linkcolor=blue!50!black,
    urlcolor=cyan!50!black,
    pdftitle={Documentation technique JobAI},
    pdfauthor={OpenAI Codex}
}

\definecolor{jobnavy}{HTML}{0F172A}
\definecolor{jobcyan}{HTML}{06B6D4}
\definecolor{jobblue}{HTML}{2563EB}
\definecolor{jobgreen}{HTML}{10B981}
\definecolor{joblight}{HTML}{F8FAFC}
\definecolor{jobgray}{HTML}{475569}
\definecolor{codebg}{HTML}{F8FAFC}
\definecolor{codeframe}{HTML}{CBD5E1}
\definecolor{codered}{HTML}{BE123C}
\definecolor{codegreen}{HTML}{15803D}
\definecolor{codepurple}{HTML}{7C3AED}

\pagestyle{fancy}
\fancyhf{}
\fancyhead[L]{JobAI}
\fancyhead[R]{Documentation technique}
\fancyfoot[C]{\thepage}
\setlength{\headheight}{14pt}

\titleformat{\section}{\Large\bfseries\color{jobnavy}}{\thesection}{0.7em}{}
\titleformat{\subsection}{\large\bfseries\color{jobblue}}{\thesubsection}{0.7em}{}
\titleformat{\subsubsection}{\normalsize\bfseries\color{jobcyan}}{\thesubsubsection}{0.7em}{}

\lstdefinestyle{jobstyle}{
    backgroundcolor=\color{codebg},
    basicstyle=\ttfamily\small,
    keywordstyle=\color{jobblue}\bfseries,
    commentstyle=\color{codegreen},
    stringstyle=\color{codered},
    numberstyle=\tiny\color{jobgray},
    numbers=left,
    stepnumber=1,
    numbersep=10pt,
    showstringspaces=false,
    breaklines=true,
    frame=single,
    rulecolor=\color{codeframe},
    tabsize=2,
    captionpos=b
}
\lstset{style=jobstyle}

\begin{document}

\begin{titlepage}
\begin{tikzpicture}[remember picture,overlay]
    \fill[jobnavy] (current page.south west) rectangle (current page.north east);
    \fill[jobblue!18] (current page.north east) circle (4cm);
    \fill[jobcyan!18] (current page.south west) circle (5cm);
    \draw[line width=2pt, color=white!20] ([xshift=1cm,yshift=-1cm]current page.north west) rectangle ([xshift=-1cm,yshift=1cm]current page.south east);
\end{tikzpicture}

\vspace*{2cm}

\begin{center}
{\Huge\bfseries\color{white} Documentation Technique Premium}\\[0.4cm]
{\LARGE\bfseries\color{white} Plateforme JobAI}\\[0.6cm]
{\Large\color{white!85} Laravel + Blade/Tailwind CSS + FastAPI + Machine Learning}\\[1.6cm]

\begin{tikzpicture}
    \node[rounded corners=18pt, fill=white, minimum width=11cm, minimum height=2.6cm] {
        \begin{minipage}{9.7cm}
        \centering
        {\Large\bfseries\color{jobnavy} Pr\'ediction de salaire et gestion d'offres d'emploi}\\[0.2cm]
        {\normalsize\color{jobgray} Rapport structur\'e avec architecture, commandes, flux m\'etier et extraits de code importants}
        \end{minipage}
    };
\end{tikzpicture}
\end{center}

\vfill

\begin{flushleft}
\renewcommand{\arraystretch}{1.5}
\begin{tabular}{>{\bfseries}p{4.2cm}p{9cm}}
Projet & JobAI \\
Architecture & Monolithe Laravel + microservice FastAPI \\
Frontend & Blade + Tailwind CSS + Chart.js \\
Backend IA & FastAPI + scikit-learn + pandas \\
Livrable & Rapport LaTeX complet \\
Date & \today \\
\end{tabular}
\end{flushleft}

\vfill
\begin{center}
{\large\color{white!80} R\'edig\'e pour documenter l'impl\'ementation r\'eelle du projet pr\'esent dans le workspace}
\end{center}
\end{titlepage}

\tableofcontents
\newpage

\section{Pr\'esentation g\'en\'erale}
JobAI est une plateforme web qui combine deux objectifs:
\begin{itemize}[leftmargin=1.2cm]
    \item estimer un salaire \`a partir d'un profil candidat;
    \item publier et consulter des offres d'emploi.
\end{itemize}

L'application repose sur une s\'eparation claire des responsabilit\'es:
\begin{itemize}[leftmargin=1.2cm]
    \item \textbf{Laravel} g\`ere l'interface, les routes web, la validation c\^ot\'e serveur et la base de donn\'ees;
    \item \textbf{Blade + Tailwind CSS} assurent le rendu frontend moderne;
    \item \textbf{FastAPI} expose le service de pr\'ediction;
    \item \textbf{scikit-learn} fournit le mod\`ele de machine learning utilis\'e pour estimer le salaire.
\end{itemize}

\section{Objectifs fonctionnels}
\begin{enumerate}[leftmargin=1.2cm]
    \item Permettre \`a un utilisateur de renseigner son exp\'erience, son dipl\^ome, son poste, sa ville et sa comp\'etence principale.
    \item Envoyer ces donn\'ees au microservice FastAPI.
    \item Recevoir une pr\'ediction salariale.
    \item Sauvegarder la pr\'ediction dans la table \texttt{predictions}.
    \item Afficher des statistiques agr\'eg\'ees par poste.
    \item G\'erer une mini-bourse d'offres d'emploi via Laravel.
\end{enumerate}

\section{Diagramme d'architecture}
\begin{figure}[H]
\centering
\begin{tikzpicture}[
    font=\small,
    layer/.style={rounded corners=14pt, thick, minimum width=13cm, minimum height=1.9cm, align=center},
    ui/.style={layer, draw=jobnavy, fill=joblight},
    app/.style={layer, draw=jobblue, fill=jobblue!8},
    service/.style={layer, draw=jobcyan!70!black, fill=jobcyan!10},
    data/.style={layer, draw=jobgreen!60!black, fill=jobgreen!10},
    flow/.style={-{Latex[length=3mm]}, thick, color=jobnavy}
]
    \node[ui] (frontend) {\textbf{Couche pr\'esentation}\\Utilisateur + interface Blade/Tailwind CSS};
    \node[app, below=0.9cm of frontend] (laravel) {\textbf{Couche applicative Laravel}\\Routes web, validation, contr\^oleurs, vues de r\'esultat et gestion des offres};
    \node[service, below=0.9cm of laravel] (fastapi) {\textbf{Couche service IA}\\API FastAPI charg\'ee d'ex\'ecuter la pr\'ediction salariale};
    \node[data, below=0.9cm of fastapi] (storage) {\textbf{Couche donn\'ees}\\MySQL (\texttt{predictions}, \texttt{job}) + mod\`ele \texttt{salary\_model.pkl} + jeu \texttt{salaries.csv}};

    \draw[flow] (frontend) -- node[right]{saisie utilisateur} (laravel);
    \draw[flow] (laravel) -- node[right]{requ\^ete HTTP JSON} (fastapi);
    \draw[flow] (fastapi) -- node[right]{lecture du mod\`ele et inf\'erence} (storage);
    \draw[flow] (fastapi.east) .. controls +(2.4,0) and +(2.4,0) .. node[right]{retour salaire estim\'e} (laravel.east);
    \draw[flow] (laravel.west) .. controls +(-2.4,0) and +(-2.4,0) .. node[left]{affichage et persistance} (frontend.west);
\end{tikzpicture}
\caption{Architecture logique en couches du projet JobAI}
\end{figure}

\subsection{Lecture du flux}
\begin{enumerate}[leftmargin=1.2cm]
    \item L'utilisateur remplit le formulaire de pr\'ediction sur l'interface Blade.
    \item Laravel valide les champs.
    \item Laravel appelle FastAPI via HTTP.
    \item FastAPI charge le mod\`ele \texttt{salary\_model.pkl} et calcule la pr\'ediction.
    \item Laravel enregistre la valeur pr\'edite en base.
    \item La vue \texttt{result.blade.php} affiche le salaire estim\'e.
\end{enumerate}

\section{Tableau des commandes}
\renewcommand{\arraystretch}{1.4}
\begin{longtable}{p{3.5cm}p{4.5cm}p{7cm}}
\toprule
\textbf{Bloc} & \textbf{Commande} & \textbf{R\^ole} \\
\midrule
\endhead
Laravel & \texttt{php artisan serve} & D\'emarre le serveur web Laravel local. \\
Laravel & \texttt{php artisan migrate} & Cr\'ee les tables de base de donn\'ees, dont \texttt{predictions}. \\
Laravel & \texttt{php artisan route:list} & Affiche les routes disponibles du projet. \\
Laravel & \texttt{npm install} & Installe les d\'ependances frontend si Vite est utilis\'e. \\
Laravel & \texttt{npm run dev} & Lance le bundler frontend en mode d\'eveloppement. \\
FastAPI & \texttt{pip install -r requirements.txt} & Installe FastAPI, pandas, scikit-learn et les biblioth\`eques Python. \\
FastAPI & \texttt{python train\_model.py} & Entra\^\i ne le mod\`ele de pr\'ediction puis g\'en\`ere \texttt{salary\_model.pkl}. \\
FastAPI & \texttt{uvicorn app:app --reload --port 8000} & D\'emarre l'API FastAPI pour la pr\'ediction. \\
Base de donn\'ees & \texttt{php artisan tinker} & Permet de tester les mod\`eles Eloquent rapidement. \\
Test API & \texttt{curl -X POST http://127.0.0.1:8000/predict-salary ...} & V\'erifie qu'une pr\'ediction JSON est bien retourn\'ee. \\
\bottomrule
\end{longtable}

\section{Structure logique du projet}
La structure du projet peut \^etre lue par blocs fonctionnels plut\^ot que sous forme de tableau, ce qui rend la logique d'ensemble plus naturelle \`a suivre.

\subsection{Noyau Laravel}
\begin{itemize}[leftmargin=1.2cm]
    \item \texttt{platforme/routes/web.php} centralise les routes web de l'application.
    \item \texttt{platforme/app/Http/Controllers/PredictionController.php} orchestre la pr\'ediction, l'enregistrement et l'historique.
    \item \texttt{platforme/app/Http/Controllers/JobController.php} pilote les actions li\'ees aux offres d'emploi.
    \item \texttt{platforme/app/Models/Prediction.php} repr\'esente les pr\'edictions stock\'ees en base.
    \item \texttt{platforme/app/Models/Job.php} mod\'elise les offres publi\'ees dans la plateforme.
\end{itemize}

\subsection{Couche interface}
\begin{itemize}[leftmargin=1.2cm]
    \item \texttt{platforme/resources/views/welcome.blade.php} contient le formulaire principal de saisie.
    \item \texttt{platforme/resources/views/result.blade.php} affiche le salaire calcul\'e.
    \item \texttt{platforme/resources/views/history.blade.php} pr\'esente l'historique et les graphiques statistiques.
    \item \texttt{platforme/resources/views/jobs.blade.php} liste les opportunit\'es disponibles.
    \item \texttt{platforme/resources/views/admin/create-job.blade.php} permet l'ajout d'une nouvelle offre.
\end{itemize}

\subsection{Moteur d'intelligence artificielle}
\begin{itemize}[leftmargin=1.2cm]
    \item \texttt{fastapi2/app.py} expose le service FastAPI de pr\'ediction.
    \item \texttt{fastapi2/train\_model.py} entra\^\i ne le mod\`ele de machine learning.
    \item \texttt{fastapi2/data/salaries.csv} fournit les donn\'ees utilis\'ees pour l'apprentissage.
    \item \texttt{fastapi2/salary\_model.pkl} correspond au mod\`ele final charg\'e lors des inf\'erences.
\end{itemize}

\section{Etapes de fonctionnement du syst\`eme}

\subsection{Etape 1: saisie du formulaire}
L'utilisateur renseigne les champs suivants:
\begin{itemize}[leftmargin=1.2cm]
    \item ann\'ees d'exp\'erience;
    \item dipl\^ome;
    \item ville;
    \item poste cibl\'e;
    \item comp\'etence principale.
\end{itemize}

\subsection{Etape 2: validation Laravel}
Le contr\^oleur Laravel s'assure que les champs sont pr\'esents et correctement typ\'es avant de contacter le moteur de pr\'ediction.

\subsection{Etape 3: appel HTTP vers FastAPI}
Laravel envoie un payload JSON au point d'entr\'ee \texttt{/predict-salary}.

\subsection{Etape 4: pr\'ediction ML}
FastAPI normalise les donn\'ees, construit un \texttt{DataFrame}, appelle \texttt{model.predict()} et renvoie le salaire estim\'e.

\subsection{Etape 5: persistance}
Le salaire estim\'e est enregistr\'e dans la table \texttt{predictions} pour alimenter l'historique.

\subsection{Etape 6: statistiques}
Laravel agr\`ege les pr\'edictions par \texttt{job\_title} et affiche la moyenne par poste.

\section{Backend Laravel}

\subsection{Routes web}
Le fichier de routes centralise la navigation entre la page d'accueil, la pr\'ediction, l'historique et les offres d'emploi.

\begin{lstlisting}[language=PHP, caption={Extrait de routes Laravel}]
use App\Http\Controllers\PredictionController;
use App\Http\Controllers\JobController;

Route::get('/', [PredictionController::class, 'index']);
Route::post('/predict', [PredictionController::class, 'predict']);
Route::get('/history', [PredictionController::class, 'history']);

Route::get('/jobs', [JobController::class, 'index']);
Route::get('/admin/jobs/create', [JobController::class, 'create']);
Route::post('/admin/jobs/store', [JobController::class, 'store']);
\end{lstlisting}

\textbf{Explication:} ces routes lient directement les URL aux contr\^oleurs. Elles constituent la colonne vert\'ebrale du backend Laravel.

\subsection{Contr\^oleur de pr\'ediction}
Ce contr\^oleur valide la requ\^ete, appelle FastAPI et sauvegarde le r\'esultat.

\begin{lstlisting}[language=PHP, caption={Extrait important de PredictionController.php}]
public function predict(Request $request)
{
    $request->validate([
        'experience_years' => 'required|numeric',
        'degree' => 'required',
        'job_title' => 'required',
        'city' => 'required',
        'skill' => 'required'
    ]);

    $response = Http::post('http://127.0.0.1:8000/predict-salary', [
        'experience_years' => (float) $request->experience_years,
        'degree' => $request->degree,
        'job_title' => $request->job_title,
        'city' => $request->city,
        'skill' => $request->skill,
    ]);

    $salary = $response['predicted_salary'];

    Prediction::create([
        'user_id' => null,
        'experience_years' => $request->experience_years,
        'degree' => $request->degree,
        'job_title' => $request->job_title,
        'city' => $request->city,
        'skill' => $request->skill,
        'predicted_salary' => $salary,
    ]);

    return view('result', compact('salary'));
}
\end{lstlisting}

\textbf{Ce que fait ce code:}
\begin{itemize}[leftmargin=1.2cm]
    \item applique une validation serveur;
    \item appelle le microservice Python;
    \item extrait la valeur \texttt{predicted\_salary};
    \item persiste les donn\'ees dans la base;
    \item affiche la page de r\'esultat.
\end{itemize}

\subsection{Historique des pr\'edictions}
L'historique exploite Eloquent pour calculer une moyenne salariale par m\'etier.

\begin{lstlisting}[language=PHP, caption={Agr\'egation des statistiques}]
$predictions = Prediction::select('job_title')
    ->selectRaw('AVG(predicted_salary) as avg_salary')
    ->groupBy('job_title')
    ->get();

$labels = $predictions->pluck('job_title');
$salaries = $predictions->pluck('avg_salary');

return view('history', compact('predictions', 'labels', 'salaries'));
\end{lstlisting}

\textbf{Int\'er\^et:} ce bloc transforme les pr\'edictions stock\'ees en indicateurs exploitables pour le tableau de bord.

\subsection{Gestion des offres d'emploi}
Le contr\^oleur suivant permet d'ajouter de nouvelles offres dans la table \texttt{job}.

\begin{lstlisting}[language=PHP, caption={Extrait de JobController.php}]
public function store(Request $request)
{
    $request->validate([
        'title' => 'required',
        'company' => 'required',
        'city' => 'required',
        'email' => 'required|email',
        'description' => 'required'
    ]);

    Job::create($request->all());

    return redirect('/jobs')->with('success', 'Job ajoute avec succes');
}
\end{lstlisting}

\subsection{Mod\`ele Eloquent Prediction}
\begin{lstlisting}[language=PHP, caption={Modele Prediction}]
class Prediction extends Model
{
    protected $fillable = [
        'user_id',
        'experience_years',
        'degree',
        'job_title',
        'city',
        'skill',
        'predicted_salary'
    ];
}
\end{lstlisting}

\textbf{Explication:} la propri\'et\'e \texttt{\$fillable} autorise l'insertion de masse en s\'ecurisant les colonnes autoris\'ees.

\subsection{Migration de la table predictions}
\begin{lstlisting}[language=PHP, caption={Migration importante pour le stockage des pr\'edictions}]
Schema::create('predictions', function (Blueprint $table) {
    $table->id();
    $table->unsignedBigInteger('user_id')->nullable();
    $table->float('experience_years');
    $table->string('degree');
    $table->string('job_title');
    $table->string('city');
    $table->string('skill');
    $table->float('predicted_salary');
    $table->timestamps();
});
\end{lstlisting}

\section{Frontend Blade et Tailwind CSS}

\subsection{Vue d'accueil et formulaire de pr\'ediction}
La vue \texttt{welcome.blade.php} sert de page d'atterrissage et de formulaire principal. Elle combine:
\begin{itemize}[leftmargin=1.2cm]
    \item une navigation claire;
    \item une mise en page responsive;
    \item un design premium avec d\'egrad\'es, cartes et boutons modernes;
    \item un formulaire structur\'e autour de champs \texttt{select} et \texttt{input}.
\end{itemize}

\begin{lstlisting}[language=HTML, caption={Extrait Blade/Tailwind du formulaire principal}]
<form method="POST" action="/predict" class="space-y-4">
    @csrf

    <div>
        <label for="experience_years"
            class="mb-2 block text-sm font-semibold text-slate-700">
            Annees d'experience
        </label>
        <input
            id="experience_years"
            type="number"
            step="0.1"
            name="experience_years"
            class="w-full rounded-2xl border border-slate-200 bg-slate-50
                   px-4 py-3 outline-none transition
                   focus:border-cyan-500 focus:bg-white
                   focus:ring-4 focus:ring-cyan-100"
            required
        >
    </div>

    <button type="submit"
        class="w-full rounded-2xl bg-slate-950 px-6 py-4
               text-sm font-semibold text-white transition hover:bg-cyan-700">
        Lancer la prediction
    </button>
</form>
\end{lstlisting}

\textbf{Pourquoi ce code est important:}
\begin{itemize}[leftmargin=1.2cm]
    \item \texttt{@csrf} prot\`ege la requ\^ete contre les attaques CSRF;
    \item les classes Tailwind acc\'el\`erent le design responsive;
    \item le formulaire envoie directement les donn\'ees \`a Laravel.
\end{itemize}

\subsection{Page de r\'esultat}
La vue \texttt{result.blade.php} valorise la pr\'ediction obtenue et met en avant la somme estim\'ee.

\begin{lstlisting}[language=HTML, caption={Bloc d'affichage du salaire estime}]
<p class="mt-4 text-5xl font-black text-slate-950 md:text-6xl">
    {{ number_format($salary, 0, ',', ' ') }}
    <span class="text-2xl font-bold text-emerald-700 md:text-3xl">DT</span>
</p>
\end{lstlisting}

\textbf{Apport UX:} le r\'esultat est lisible, central, bien hi\'erarchis\'e et coh\'erent avec une interface premium.

\subsection{Tableau de bord statistique}
La vue \texttt{history.blade.php} utilise Chart.js pour transformer des donn\'ees PHP en graphique interactif.

\begin{lstlisting}[language=JavaScript, caption={Injection des donnees PHP vers Chart.js}]
const labels = @json($labels);
const data = @json($salaries);

new Chart(ctx, {
    type: 'bar',
    data: {
        labels: labels,
        datasets: [{
            label: 'Salaire moyen (DT)',
            data: data
        }]
    }
});
\end{lstlisting}

\textbf{Id\'ee cl\'e:} Laravel pr\'epare les donn\'ees, Blade les injecte dans le JavaScript, puis Chart.js les visualise.

\subsection{Vue offres d'emploi}
La vue \texttt{jobs.blade.php} affiche les annonces sous forme de cartes modernes, tandis que \texttt{create-job.blade.php} permet l'ajout d'une nouvelle offre.

\begin{lstlisting}[language=HTML, caption={Extrait Blade d'une carte offre d'emploi}]
@foreach($jobs as $job)
    <article class="group flex h-full flex-col rounded-[2rem] border
                    border-white/70 bg-white p-6 shadow-sm">
        <h2 class="mt-1 text-xl font-bold text-slate-950">{{ $job->title }}</h2>
        <p class="mt-1 text-sm text-slate-500">{{ $job->company }}</p>
        <span class="rounded-full bg-slate-100 px-3 py-1
                     text-xs font-semibold text-slate-700">{{ $job->city }}</span>
    </article>
@endforeach
\end{lstlisting}

\section{FastAPI et mod\`ele de pr\'ediction}

\subsection{Chargement du mod\`ele et des donn\'ees}
Le microservice FastAPI charge au d\'emarrage le mod\`ele entrain\'e et le fichier CSV contenant les valeurs de r\'ef\'erence.

\begin{lstlisting}[language=Python, caption={Chargement du modele et des options de formulaire}]
MODEL_PATH = Path("salary_model.pkl")
DATA_PATH = Path("data/salaries.csv")

with open(MODEL_PATH, "rb") as f:
    model = pickle.load(f)

dataset = pd.read_csv(DATA_PATH)

FORM_OPTIONS = {
    "degree": sorted(dataset["degree"].dropna().astype(str).unique().tolist()),
    "job_title": sorted(dataset["job_title"].dropna().astype(str).unique().tolist()),
    "city": sorted(dataset["city"].dropna().astype(str).unique().tolist()),
    "skill": sorted(dataset["skill"].dropna().astype(str).unique().tolist()),
}
\end{lstlisting}

\textbf{R\^ole:} ce bloc rend l'API capable de valider les valeurs attendues et d'offrir une coh\'erence m\'etier entre les donn\'ees d'entra\^\i nement et l'inf\'erence.

\subsection{Sch\'ema de requ\^ete}
\begin{lstlisting}[language=Python, caption={Schema Pydantic pour la prediction}]
class SalaryPredictionRequest(BaseModel):
    experience_years: float
    degree: str
    job_title: str
    city: str
    skill: str
\end{lstlisting}

\textbf{Explication:} Pydantic structure les entr\'ees, simplifie la validation et documente implicitement l'API.

\subsection{Validation et normalisation}
\begin{lstlisting}[language=Python, caption={Fonctions de preparation des donnees}]
def normalize(payload: dict):
    return {
        "experience_years": payload["experience_years"],
        "degree": payload["degree"].strip(),
        "job_title": payload["job_title"].strip(),
        "city": payload["city"].strip(),
        "skill": payload["skill"].strip(),
    }

def validate(payload: dict):
    errors = []

    if payload["experience_years"] < 0 or payload["experience_years"] > 50:
        errors.append("experience_years must be between 0 and 50")

    for f in ["degree", "job_title", "city", "skill"]:
        if not payload[f]:
            errors.append(f"{f} is required")

    for f in ["degree", "job_title", "city", "skill"]:
        if payload[f] not in FORM_OPTIONS[f]:
            errors.append(f"{f} invalid value")

    return errors
\end{lstlisting}

\textbf{Int\'er\^et technique:}
\begin{itemize}[leftmargin=1.2cm]
    \item nettoyage des cha\^\i nes;
    \item contr\^ole des bornes sur l'exp\'erience;
    \item v\'erification des valeurs autoris\'ees.
\end{itemize}

\subsection{Endpoint principal de pr\'ediction}
\begin{lstlisting}[language=Python, caption={Endpoint FastAPI pour Laravel}]
@app.post("/predict-salary")
def predict_simple(data: SalaryPredictionRequest):
    payload = normalize(data.model_dump())
    prediction = predict_salary(payload)

    return {
        "predicted_salary": prediction
    }
\end{lstlisting}

\textbf{Lecture:} Laravel appelle cet endpoint, re\c{c}oit un JSON simple et peut afficher le salaire sans transformation complexe.

\subsection{Fonction d'inf\'erence}
\begin{lstlisting}[language=Python, caption={Inference scikit-learn}]
def predict_salary(payload: dict):
    df = pd.DataFrame([payload])
    pred = model.predict(df)[0]
    return round(float(pred), 2)
\end{lstlisting}

\textbf{Point cl\'e:} la pr\'ediction repose sur une seule ligne de donn\'ees transform\'ee en \texttt{DataFrame} pour rester compatible avec le pipeline scikit-learn.

\section{Script d'entra\^\i nement du mod\`ele}

\subsection{Construction du pipeline ML}
Le script \texttt{train\_model.py} cr\'ee un pipeline comprenant:
\begin{itemize}[leftmargin=1.2cm]
    \item \texttt{OneHotEncoder} pour les variables cat\'egorielles;
    \item \texttt{passthrough} pour la variable num\'erique \texttt{experience\_years};
    \item \texttt{RandomForestRegressor} pour la r\'egression finale.
\end{itemize}

\begin{lstlisting}[language=Python, caption={Pipeline de machine learning}]
def build_pipeline() -> Pipeline:
    categorical_features = ["degree", "job_title", "city", "skill"]
    numeric_features = ["experience_years"]

    preprocessor = ColumnTransformer(
        transformers=[
            ("cat", OneHotEncoder(handle_unknown="ignore"), categorical_features),
            ("num", "passthrough", numeric_features),
        ]
    )

    model = RandomForestRegressor(
        n_estimators=300,
        max_depth=12,
        min_samples_split=2,
        min_samples_leaf=1,
        random_state=42,
    )

    return Pipeline([("preprocessor", preprocessor), ("model", model)])
\end{lstlisting}

\subsection{Entra\^\i nement et sauvegarde}
\begin{lstlisting}[language=Python, caption={Apprentissage, evaluation et export du modele}]
df = load_dataset()
X = df[FEATURES]
y = df["salary"]
pipeline = build_pipeline()

X_train, X_test, y_train, y_test = train_test_split(
    X, y, test_size=0.2, random_state=42
)

pipeline.fit(X_train, y_train)
predictions = pipeline.predict(X_test)

mae = mean_absolute_error(y_test, predictions)
r2 = r2_score(y_test, predictions)

with open(MODEL_PATH, "wb") as file_handle:
    pickle.dump(pipeline, file_handle)
\end{lstlisting}

\textbf{Analyse:}
\begin{itemize}[leftmargin=1.2cm]
    \item \texttt{MAE} mesure l'erreur absolue moyenne;
    \item \texttt{R2} mesure la qualit\'e d'ajustement du mod\`ele;
    \item le pipeline entier est s\'erialis\'e dans \texttt{salary\_model.pkl}.
\end{itemize}

\section{Contrat d'\'echange entre Laravel et FastAPI}
\subsection{Requ\^ete envoy\'ee}
\begin{lstlisting}[language=json, caption={Exemple de payload JSON}]
{
  "experience_years": 3.5,
  "degree": "Master",
  "job_title": "Data Scientist",
  "city": "Tunis",
  "skill": "Python"
}
\end{lstlisting}

\subsection{R\'eponse retourn\'ee}
\begin{lstlisting}[language=json, caption={Reponse JSON simplifiee}]
{
  "predicted_salary": 3200.75
}
\end{lstlisting}

\section{Synth\`ese technique par couche}
\begin{longtable}{p{3.5cm}p{4.5cm}p{6.5cm}}
\toprule
\textbf{Couche} & \textbf{Technologie} & \textbf{Responsabilit\'e} \\
\midrule
\endhead
Interface & Blade + Tailwind CSS & Affichage, formulaire, UX premium, responsive design. \\
Contr\^ole m\'etier & Laravel Controllers & Validation, orchestration du flux, persistance. \\
Persistance & Eloquent + migrations & Stockage des pr\'edictions et des offres. \\
Visualisation & Chart.js & Repr\'esentation graphique des moyennes salariales. \\
IA / API & FastAPI + Pydantic & Exposition du service de pr\'ediction. \\
Machine Learning & pandas + scikit-learn & Entra\^\i nement, encodage, pr\'ediction du salaire. \\
\bottomrule
\end{longtable}

\section{Points forts du projet}
\begin{itemize}[leftmargin=1.2cm]
    \item architecture simple \`a comprendre et facile \`a d\'emontrer;
    \item s\'eparation nette entre interface web et moteur de pr\'ediction;
    \item exploitation d'un vrai pipeline machine learning;
    \item design frontend moderne avec Tailwind CSS;
    \item exploitation des pr\'edictions stock\'ees pour produire des statistiques.
\end{itemize}

\section{Conclusion}
Le projet JobAI illustre une architecture hybride tr\`es pertinente pour un projet acad\'emique ou professionnel: Laravel pilote l'exp\'erience utilisateur et la logique web, tandis que FastAPI fournit une brique IA sp\'ecialis\'ee pour la pr\'ediction de salaire. Le frontend Blade/Tailwind CSS apporte une pr\'esentation moderne, et la couche machine learning permet de valoriser les donn\'ees par une estimation concr\`ete et exploitable.

Ce document LaTeX regroupe les \'el\'ements les plus importants du projet: page de garde premium, diagramme d'architecture, tableau des commandes, extraits de code essentiels et explications claires des diff\'erentes \'etapes.

\end{document}
